\chapter{Matriz de triangulación (encuesta y entrevistas)}
\label{apx:triangulacion}

\section*{1. Criterio de integración}
Se integraron los resultados de la encuesta a estudiantes con las perspectivas obtenidas mediante las entrevistas semiestructuradas. Ambas técnicas aportan evidencia de distinta naturaleza: la encuesta describe percepciones de los 213 estudiantes que integraron la muestra, mientras que las entrevistas aportan perspectivas individuales vinculadas con la docencia, la investigación y la gestión académica. En consecuencia, no se trataron como mediciones equivalentes del mismo fenómeno.

La triangulación se realizó para organizar relaciones de convergencia, complementariedad, diferencias y evidencia aportada por una única fuente, conservando el alcance propio de cada técnica. La matriz que se presenta a continuación reproduce las once dimensiones de la triangulación tal como se utilizaron en el análisis, con los mismos cálculos de la encuesta y la misma codificación de las entrevistas.

\section*{2. Matriz de triangulación}
\begin{center}\footnotesize
\begin{infobox}
\textbf{\textcolor{accent}{1. Confianza e integridad en evaluaciones digitales}} \textit{(Nivel de respaldo: Fuerte)}\\
\textit{Evidencia de encuesta:} EST05: media 2,68; 40,4\,\% en 1--2, 39,9\,\% en 3 y solo 19,7\,\% en 4--5 de confiabilidad.\\
\textit{Evidencia de entrevistas:} ENT-01 plantea dificultades para controlar fehacientemente una evaluación remota; ENT-02 desplaza la validez hacia el diseño evaluativo; ENT-03 relaciona integridad con alineación entre objetivos, enseñanza, instrumento y recursos permitidos.\\
\textit{Relación entre fuentes:} Complementariedad con convergencia conceptual. No miden lo mismo: estudiantes expresan confiabilidad; entrevistados explican condiciones de validez/integridad.\\
\textit{Interpretación:} La confiabilidad de la evaluación digital aparece como un problema relevante, pero no puede atribuirse a una única causa tecnológica.\\
\textit{Implicación preliminar:} Abordar la integridad como problema de trazabilidad, diseño y procedimiento, no solamente de detección de copia.
\end{infobox}
\vspace{2pt}
\begin{infobox}
\textbf{\textcolor{accent}{2. Multiplicidad y contexto de los riesgos}} \textit{(Nivel de respaldo: Fuerte)}\\
\textit{Evidencia de encuesta:} EST07: internet 80,3\,\%, materiales no permitidos 76,1\,\%, comunicación entre estudiantes 66,7\,\%, ayuda externa 64,8\,\%; además 58,2\,\% identificó problemas técnicos que afectan el control.\\
\textit{Evidencia de entrevistas:} Las entrevistas mencionan asistencia externa, IA, materiales, identidad, entorno, dispositivos y condiciones técnicas. ENT-02 y ENT-03 destacan que el significado de utilizar un recurso depende del objetivo de la evaluación.\\
\textit{Relación entre fuentes:} Convergencia + complementariedad.\\
\textit{Interpretación:} El problema no admite una definición reducida a una única conducta ni parece razonable depender de una señal aislada.\\
\textit{Implicación preliminar:} Considerar múltiples señales y mantener contexto académico para interpretarlas.
\end{infobox}
\vspace{2pt}
\begin{infobox}
\textbf{\textcolor{accent}{3. Invasividad y privacidad}} \textit{(Nivel de respaldo: Fuerte para la relativización del supuesto; moderado para comparaciones entre tecnologías.)}\\
\textit{Evidencia de encuesta:} Cámara: media 2,34 y 56,3\,\% en 1--2 de invasividad. Pantalla: media 2,98. Registro de eventos: media 3,09, con 40,8\,\% en 4--5. COM03 muestra aceptación heterogénea: 37,1\,\% alta, 36,2\,\% intermedia, 26,8\,\% baja.\\
\textit{Evidencia de entrevistas:} ENT-01 relata preocupaciones por protección de datos y control del dispositivo. ENT-02 señala dificultades asociadas a cámaras, múltiples dispositivos y condiciones del entorno, y la necesidad de informar qué se registra. No existe evidencia equivalente clara en ENT-03.\\
\textit{Relación entre fuentes:} Complementariedad. Las entrevistas muestran costos de mecanismos intensivos, mientras la encuesta impide asumir que registrar eventos sea automáticamente percibido como menos invasivo.\\
\textit{Interpretación:} La privacidad no se resuelve simplemente eliminando la cámara. Importan qué se captura, con qué finalidad, cuánto se captura y cómo se comunica.\\
\textit{Implicación preliminar:} Evitar posicionar el enfoque únicamente como ``menos invasivo''; priorizar proporcionalidad, minimización y transparencia.
\end{infobox}
\vspace{2pt}
\begin{infobox}
\textbf{\textcolor{accent}{4. Transparencia previa}} \textit{(Nivel de respaldo: Moderado)}\\
\textit{Evidencia de encuesta:} COM01: 93,0\,\% asignó 4--5 a conocer previamente qué datos se registran y con qué finalidad.\\
\textit{Evidencia de entrevistas:} ENT-02 afirma que la recolección debe avisarse previamente y que reglas y consecuencias deben estar definidas de antemano. ENT-03 aporta sobre necesidad de procedimientos claros, pero ya en contexto de \textit{feedback}.\\
\textit{Relación entre fuentes:} Convergencia, con distinta calidad temporal de la evidencia cualitativa.\\
\textit{Interpretación:} La transparencia aparece como condición de legitimidad del procedimiento, no como detalle de interfaz.\\
\textit{Implicación preliminar:} Informar previamente datos, finalidad, reglas y posibles usos de las alertas.
\end{infobox}
\vspace{2pt}
\begin{infobox}
\textbf{\textcolor{accent}{5. Revisión humana}} \textit{(Nivel de respaldo: Fuerte)}\\
\textit{Evidencia de encuesta:} COM02: 93,9\,\% en 4--5. EST14: revisión humana posterior fue la alternativa individual más elegida, 51,2\,\%. COM04: revisión humana fue la condición individual más seleccionada, 30,5\,\%.\\
\textit{Evidencia de entrevistas:} ENT-01 afirma espontáneamente que ``la curaduría humana del resultado tiene que estar siempre''. ENT-02 describe decisiones caso a caso y validaciones. ENT-03 rechaza delegar autónomamente la decisión académica en IA.\\
\textit{Relación entre fuentes:} Convergencia clara.\\
\textit{Interpretación:} La interpretación de una señal con consecuencias académicas requiere conservar intervención humana.\\
\textit{Implicación preliminar:} Mantener revisión humana de señales y decisiones académicas.
\end{infobox}
\vspace{2pt}
\begin{infobox}
\textbf{\textcolor{accent}{6. Límites a acusaciones o sanciones automáticas}} \textit{(Nivel de respaldo: Fuerte)}\\
\textit{Evidencia de encuesta:} EST15: 75,1\,\% asignó valores 4--5 a la importancia de que una alerta no implique automáticamente una acusación. EST16: 80,3\,\% seleccionó evitar decisiones automáticas de sanción. EST14: anulación automática 5,2\,\%.\\
\textit{Evidencia de entrevistas:} ENT-01 cuestiona mecanismos de expulsión automática por contingencias técnicas. ENT-03 distingue automatizar reglas parametrizadas de permitir que la IA decida académicamente. ENT-02 mantiene el criterio docente/contextual.\\
\textit{Relación entre fuentes:} Convergencia.\\
\textit{Interpretación:} Automatizar detección o tareas operativas no equivale a automatizar juicio académico o sanción.\\
\textit{Implicación preliminar:} Separar generación de indicadores de acusación, invalidación o sanción.
\end{infobox}
\vspace{2pt}
\end{center}
\begin{center}\footnotesize
\begin{infobox}
\textbf{\textcolor{accent}{7. Evidencia suficiente y accionabilidad de las alertas}} \textit{(Nivel de respaldo: Moderado)}\\
\textit{Evidencia de encuesta:} La encuesta no mide directamente ``suficiencia probatoria''. Sí muestra preferencia por revisión humana y por una instancia adicional de validación: 20,7\,\% en EST14.\\
\textit{Evidencia de entrevistas:} ENT-01 plantea que un cambio de pestaña no necesariamente permite sostener una sanción; ENT-02 propone usar señales para seleccionar casos a validar; ENT-03 pregunta explícitamente para qué sirve obtener información si luego no habilita una acción clara.\\
\textit{Relación entre fuentes:} Complementariedad, no convergencia métrica.\\
\textit{Interpretación:} Una alerta puede ser útil para decidir investigar más sin ser suficiente para una consecuencia disciplinaria. El estándar de evidencia debe depender de la acción posterior.\\
\textit{Implicación preliminar:} Diseñar evidencia para orientar decisiones proporcionales y acciones posteriores claras.
\end{infobox}
\vspace{2pt}
\begin{infobox}
\textbf{\textcolor{accent}{8. Falsos positivos y ambigüedad de señales}} \textit{(Nivel de respaldo: Moderado)}\\
\textit{Evidencia de encuesta:} No existe una medición cuantitativa directa de falsos positivos. Las preferencias por revisión y no automatización son evidencia indirecta, no equivalente.\\
\textit{Evidencia de entrevistas:} ENT-01 describe incidencias que deben distinguirse de falsos positivos. ENT-02 relata alertas producidas por movimientos naturales de la mirada y presión adicional pese a no estar haciendo trampa. Ambos aparecen en \textit{discovery}.\\
\textit{Relación entre fuentes:} Evidencia directa solo en entrevistas; encuesta complementaria indirecta.\\
\textit{Interpretación:} Un evento observable admite explicaciones alternativas. La anomalía no debería tratarse como evidencia autosuficiente de intención o culpabilidad.\\
\textit{Implicación preliminar:} Presentar señales con contexto y permitir su interpretación, evitando conclusiones binarias automáticas.
\end{infobox}
\vspace{2pt}
\begin{infobox}
\textbf{\textcolor{accent}{9. Validaciones complementarias}} \textit{(Nivel de respaldo: Moderado)}\\
\textit{Evidencia de encuesta:} EST14: 20,7\,\% eligió explícitamente una instancia adicional de validación, segunda opción más frecuente después de revisión humana.\\
\textit{Evidencia de entrevistas:} ENT-02 propone oral focalizado; ENT-03 repreguntas o preguntas de validación; ENT-01 describe el oral como mecanismo utilizado previamente, aunque advierte que no es suficiente como único control.\\
\textit{Relación entre fuentes:} Convergencia parcial.\\
\textit{Interpretación:} Una alerta incierta puede servir para focalizar una verificación adicional del conocimiento en vez de transformarse directamente en sanción.\\
\textit{Implicación preliminar:} Contemplar mecanismos de profundización o validación cuando la evidencia no sea concluyente.
\end{infobox}
\vspace{2pt}
\begin{infobox}
\textbf{\textcolor{accent}{10. Carga docente y escalabilidad}} \textit{(Nivel de respaldo: Exploratorio)}\\
\textit{Evidencia de encuesta:} No medido por la encuesta, porque la fuente cuantitativa corresponde exclusivamente a estudiantes.\\
\textit{Evidencia de entrevistas:} ENT-01 señala capacitación y revisión de reportes; ENT-02 problematiza qué hacer si 80 de 100 alumnos quedan señalados; ENT-03 cuestiona gastar recursos en información sin uso posterior.\\
\textit{Relación entre fuentes:} Evidencia disponible solo en entrevistas.\\
\textit{Interpretación:} La utilidad de una herramienta no depende únicamente de detectar eventos, sino de que el volumen de revisión posterior resulte manejable.\\
\textit{Implicación preliminar:} Reducir ruido y priorizar información para evitar una carga de revisión inviable.
\end{infobox}
\vspace{2pt}
\begin{infobox}
\textbf{\textcolor{accent}{11. Dependencia del tipo y contexto de la evaluación}} \textit{(Nivel de respaldo: Fuerte)}\\
\textit{Evidencia de encuesta:} EST06: 45,1\,\% considera más difícil garantizar condiciones equitativas digitalmente y 41,3\,\% responde ``depende del tipo de evaluación''; solo 6,6\,\% responde ``no''.\\
\textit{Evidencia de entrevistas:} ENT-02 distingue nivel de carrera, tipo de consigna y uso de IA; ENT-03 vincula herramientas permitidas con objetivos de aprendizaje; ENT-01 también diferencia usos y modalidades de control.\\
\textit{Relación entre fuentes:} Convergencia fuerte.\\
\textit{Interpretación:} No existe evidencia para sostener un único esquema de supervisión válido para todas las evaluaciones.\\
\textit{Implicación preliminar:} Ajustar la interpretación y el nivel de supervisión al objetivo, formato y reglas de cada evaluación.
\end{infobox}
\vspace{2pt}
\end{center}
Fuente: elaboración propia a partir de la encuesta y las entrevistas del \textit{user research}.

\section*{3. Nota metodológica de lectura}
La convergencia entre fuentes no significa que la encuesta y las entrevistas hayan medido exactamente lo mismo. La encuesta corresponde a percepciones autoinformadas de estudiantes y no mide conducta real ni frecuencia de fraude; las entrevistas corresponden a perspectivas individuales situadas y no se atribuyen a UADE como institución.

Cuando una dimensión cuenta con evidencia de una sola fuente, se conserva ese alcance. Del mismo modo, una señal atípica no se interpreta por sí sola como prueba de deshonestidad y la evidencia obtenida durante \textit{feedback} posterior a la presentación de ImprontIA no se presenta como \textit{discovery}. La matriz tampoco permite concluir que la ausencia de cámara implique, por sí misma, menor invasividad.

Los niveles de respaldo expresan el grado de apoyo de cada dimensión dentro de este estudio y no constituyen una generalización estadística. La triangulación orientó la formulación de los seis hallazgos principales del cuerpo del capítulo; la carga docente y la escalabilidad se conservan como una consideración cualitativa exploratoria y no como un séptimo hallazgo principal. En ningún caso la triangulación constituye una validación directa de ImprontIA.
